% DRAFT replacement for main_paper11.tex.
% Scope: (A) the proof of Proposition P11:prop:R_equals_1 (currently the
% sigma_Slater + sigma_xc = 6.00 decomposition), (B) the remark
% P11:rem:sigma_xc_origin, (C) the one-line estimate at line ~760,
% (D) the summary restatement at lines ~1882-1886.
% Not applied. Numbers verified against mendeleev (Clementi et al. and
% Slater tables) and by direct computation; see
% papers/src_paper11/zeff_prescription/.

% ============================================================== (A)
% Replaces everything between \begin{proof} and \qed \end{proof} of
% Proposition P11:prop:R_equals_1.

\begin{proof}
We work in atomic units~\eqref{P11:eq:HZ_au}; ``Ry'' below means $2\ERy$.
By Koopmans' theorem, $\IE(\mathrm{O}) = -\varepsilon_{2p}$, and writing
the orbital energy in hydrogenic form defines an effective charge,
\begin{equation}
  -\varepsilon_{2p}
  \;=\; \frac{Z_\mathrm{eff}^2}{n_\mathrm{qn}^2}\,\ERy,
  \qquad
  Z_\mathrm{eff} \;\equiv\; Z - \sigma_\mathrm{tot},
  \label{P11:eq:eps_Zeff}
\end{equation}
so that $R(\mathrm{O})=1$ is exactly $Z_\mathrm{eff}=n_\mathrm{qn}=2$.
Equation~\eqref{P11:eq:eps_Zeff} is a change of variables rather than a
derivation: it \emph{defines} $Z_\mathrm{eff}$ from the energy, so
``$Z_\mathrm{eff}(\mathrm{O})=2$'' and ``$\IE(\mathrm{O})=\ERy$'' are
the same statement.  Establishing~\eqref{P11:eq:R_equals_1} requires
$\sigma_\mathrm{tot}$ from a prescription independent of the measured
energy.

\textbf{(a) Independent screening prescriptions.}
Two are standard.  Slater's rules~\cite{Slater1930} place the $2s$ and
$2p$ electrons in a single screening group, each screening the others
at $0.35$, with the $1s$ pair at $0.85$:
\begin{equation}
  \sigma_\mathrm{Slater}
  = 2\times0.85\;(1s^2) + 5\times0.35\;(\text{other } n{=}2)
  = 3.45,
  \qquad
  Z_\mathrm{eff}^\mathrm{Slater} = 4.55 .
  \label{P11:eq:sigma_slater}
\end{equation}
The self-consistent Hartree--Fock screening constants of Clementi
\emph{et al.}~\cite{Clementi1963} give, for oxygen,
$Z_\mathrm{eff}(2s)=4.4916$ and $Z_\mathrm{eff}(2p)=4.4532$.  The two
orbitals differ by $0.038$, confirming Slater's grouping directly: $2s$
and $2p$ see essentially the same effective charge, and $2s$ does not
act as an inner shell towards $2p$.  Taking the $1s$ pair at $0.85$
each, the Clementi value implies $0.369$ of screening per $n=2$
electron against Slater's $0.35$, a five per cent refinement.  The two
prescriptions agree on $Z_\mathrm{eff}(2p)$ to two per cent.

\textbf{(b) Neither prescription gives $Z_\mathrm{eff}=2$, and none can.}
Both place $Z_\mathrm{eff}(\mathrm{O},2p)$ near $4.5$.  The value
$\sigma_\mathrm{tot}=6.00$ required by~\eqref{P11:eq:eps_Zeff} lies
outside the range any screening decomposition can reach: with the $1s$
pair at $0.85$ and $2p$--$2p$ at $0.35$, each $2s$ electron would have
to screen $1.625$ units of nuclear charge, and a single electron cannot
screen more than one; granting the $1s$ pair the maximum $1.0$ each
still requires $1.475$.  Consequently $Z_\mathrm{eff}=2$ is not an
orbital screening constant but the measured ionisation energy
re-expressed through~\eqref{P11:eq:eps_Zeff}.

\textbf{(c) The correct use of screening constants.}
Equation~\eqref{P11:eq:eps_Zeff} is not how screening constants enter
energies.  Slater's rules furnish total energies
\begin{equation}
  E \;=\; -\ERy \sum_i \Bigl(\frac{Z_{\mathrm{eff},i}}{n^*_i}\Bigr)^{\!2},
  \label{P11:eq:slater_total}
\end{equation}
with $n^*$ the effective principal quantum numbers, and the ionisation
energy is the difference between the ion and the neutral atom.
Evaluated for $\mathrm{O}\,[1s^22s^22p^4]$ and
$\mathrm{O}^+[1s^22s^22p^3]$, equation~\eqref{P11:eq:slater_total}
gives
\begin{equation}
  \IE_\mathrm{Slater}(\mathrm{O})
  \;=\; E(\mathrm{O}^+) - E(\mathrm{O})
  \;=\; 14.17\,\mathrm{eV}
  \;=\; 1.041\,\ERy,
  \qquad
  |\Delta n| = 0.041 .
  \label{P11:eq:IE_slater_total}
\end{equation}
Standard screening theory therefore places oxygen's first ionisation
energy at the Rydberg energy to four per cent, with no free parameters
and no framework input, establishing~\eqref{P11:eq:R_equals_1} at that
order.  The $O(\aem)$ refinement of
Proposition~\ref{P11:prop:Zeff_closed} is a separate, smaller effect.
\qed
\end{proof}

% ============================================================== (B)
% Replaces the remark P11:rem:sigma_xc_origin in its entirety.

\begin{remark}[Precision of the H--O coincidence]
\label{P11:rem:sigma_xc_origin}
Equation~\eqref{P11:eq:IE_slater_total} accounts for
$R(\mathrm{O})=1$ to four per cent, $|\Delta n| = 0.041$.  The measured
coincidence is far tighter: $\IE(\mathrm{H})=13.598\,\mathrm{eV}$ and
$\IE(\mathrm{O})=13.618\,\mathrm{eV}$ differ by $0.15\%$,
$|\Delta n| = 0.0015$, some twenty-seven times closer than the
screening calculation predicts.  That excess precision is not derived
here, and is the substance of the near-degeneracy.

Two accounts of it are unavailable.  A decomposition
$\sigma_\mathrm{tot} = \sigma_\mathrm{Slater} + \sigma_\mathrm{xc}$
cannot supply it, because $\sigma_\mathrm{xc}$ would be defined as the
residual $\sigma_\mathrm{tot}-\sigma_\mathrm{Slater}$, making its
closure an identity rather than a check --- and, by the ceiling
argument above, no admissible $\sigma_\mathrm{xc}$ reaches
$\sigma_\mathrm{tot}=6.00$ in any case.  Nor does the $2p^4$ exchange
penalty account for the gap quantitatively: with the Condon--Shortley
angular coefficient $3/25$ and $F^2(2p)=(17/81)Z_\mathrm{eff}\ERy$ at
$Z_\mathrm{eff}=2$, the exchange term is
$(3/2)K_{2p,2p} = 1.03\,\mathrm{eV}$, which at
$\mathrm{d}\IE/\mathrm{d}\sigma = 2Z_\mathrm{eff}\ERy/n_\mathrm{qn}^2
= 13.6\,\mathrm{eV}$ per unit of screening corresponds to
$\Delta\sigma \approx 0.08$.

What the framework contributes is unaffected: the Rydberg energy sits
at tower level $-12.102$ with no free parameters once $\aem$ and $m_e$
are fixed (Corollary~\ref{P11:cor:rydberg}), and hydrogen sits there by
definition.  That oxygen sits within $0.0015$ of the same level is an
atomic-structure fact, approximately explained by ordinary screening
and exactly explained by neither.
\end{remark}

% ============================================================== (C)
% Replaces the parenthetical estimate near line 760.

to the observed $13.618\,\mathrm{eV} \approx \ERy$.
% (delete the clause ``from the screened-Coulomb estimate of
%  $\approx 3.55^2\,\ERy/4 \approx 42.9\,\mathrm{eV}$'': the
%  one-electron orbital formula is not the right expression, and with
%  the corrected Slater charge it would read $4.55^2\ERy/4 = 70.4$ eV.
%  If a comparison number is wanted, use the total-energy value
%  $14.17\,\mathrm{eV}$ of \eqref{P11:eq:IE_slater_total}.)

% ============================================================== (D)
% Replaces the summary restatement at lines ~1882-1886.

  Slater's rules give $Z_\mathrm{eff}(\mathrm{O},2p)=4.55$ and the
  Clementi \emph{et al.} Hartree--Fock constants give $4.4532$,
  agreeing to two per cent; used as total-energy differences they
  place $\IE(\mathrm{O})$ within four per cent of $\ERy$.  The
  parametrisation $Z_\mathrm{eff}=2=Z-\sigma_\mathrm{tot}$ is
  equivalent to $\IE(\mathrm{O})=\ERy$ itself and is not an
  independent screening result.
